% ============================================================
%  Arsenic Colorimeter: Synthesis and Planning Document
%
%  Compile with: pdflatex colorimeter_plan.tex (twice)
% ============================================================

\documentclass[11pt,letterpaper]{article}

\usepackage[letterpaper,margin=1in]{geometry}
\usepackage{amsmath,amssymb}
\usepackage[protrusion=true,expansion=false]{microtype}
\usepackage{parskip}
\usepackage{enumitem}
\usepackage{hyperref}
\hypersetup{
  colorlinks=true,
  linkcolor=blue,
  urlcolor=blue,
  pdftitle={Arsenic Colorimeter -- Synthesis and Planning},
  pdfauthor={Design Notes}
}
\usepackage{titlesec}
\titleformat{\section}{\Large\bfseries}{\thesection}{1em}{}
\titleformat{\subsection}{\large\bfseries}{\thesubsection}{1em}{}
\titleformat{\subsubsection}{\normalsize\bfseries}{\thesubsubsection}{1em}{}

\setlist[itemize]{topsep=2pt,itemsep=1pt}
\setlist[enumerate]{topsep=2pt,itemsep=1pt}

\title{Synthesis and Planning:\\ Arsenic Colorimeter Methods}
\author{Design notes}
\date{\today}

\begin{document}
\maketitle

% ============================================================
\section*{Where we are}
% ============================================================

You have a hardware-first project with a working optical prototype, a clear
measurement goal (arsenic in drinking water, with the decision threshold
around the WHO 10\,\textmu g/L limit), a chemistry you don't fully control
(gas-phase Gutzeit-type strip reaction), and a statistics pipeline that was
written somewhat ahead of the hardware and the data. The original \texttt{.tex}
is a speculative methods document --- a plausible skeleton assembled from
standard tools (Mahalanobis, inverse-Wishart, GP), not a reflection of
measurements you've actually looked at.

My critique flagged two kinds of issues: mathematical consistency problems
inside the skeleton, and physical difficulties the skeleton didn't
acknowledge. Your response splits those cleanly:

\begin{itemize}
  \item \textbf{Physical issues you've already thought about and have a plan for:}
    illumination non-uniformity (polynomial flat-field), strip heterogeneity
    (defer, study empirically), strip kinetics (built-in timer, fixed
    protocol), temperature (ambient + water + LED-adjacent probes), stray
    light (flat-black interior, small patch area, testable).
  \item \textbf{Physical issues still open:} batch-to-batch strip variation,
    the exact spectral behavior of the IMX219 under each LED, HDR / dynamic
    range in the pale-yellow regime, dark current behavior vs.\ exposure
    time and temperature.
  \item \textbf{Statistical issues you're explicitly deferring to
    collaboration:} what to actually use as features, what to regress with,
    what to report for uncertainty, and how to validate.
\end{itemize}

My job at this planning stage is to (a)~design a statistical pipeline that
fits your hardware plan, (b)~tell you what data I need to see before
committing to specifics, and (c)~flag what you should learn so you're not
flying blind on the statistics.

% ============================================================
\section*{The approach I'd recommend, at the level of strategy}
% ============================================================

The original paper's architecture was: \emph{within-image Mahalanobis
distances to reference patches, stacked with LED temperatures, fed to a GP.}
I don't think that's the right shape. Here is the shape I'd recommend
instead, with the reasoning.

\textbf{Build the pipeline in layers, each with a clear job, each testable on
its own.} The layers are: raw capture $\to$ pixel-level corrections $\to$
patch-level feature extraction $\to$ dimensional reduction $\to$
regression/classification $\to$ uncertainty calibration. Each layer should
be swappable. The temptation with a project like this is to combine layers
--- e.g., fold flat-fielding into the covariance model --- and that makes
debugging impossible.

\subsection*{Layer 1: Raw capture and pixel-level corrections}

This is the part where your hardware plan already lives. Concretely:

\begin{itemize}
  \item \textbf{Dark subtraction first, always.} Every image, raw, minus a
    same-exposure-time dark frame. Dark-subtract before anything else. The
    dark frame captures read noise bias, dark current, and fixed-pattern
    noise in one shot. Your plan mentions ``subtract the mean black
    background'' but a \emph{mean} dark throws away pattern noise. Subtract
    the full dark \emph{frame}, pixel-by-pixel. Capture darks periodically
    (they drift with sensor temperature).
  \item \textbf{Flat-field correction using your polynomial plan.} Your
    plan is sound. One refinement: fit the polynomial \emph{per Bayer
    channel} (you already said this) and \emph{per LED} (you already said
    this), but also consider fitting across a wider area than just the
    rectangles and evaluating at rectangle centers. A polynomial fit over
    more pixels is better conditioned. Degree-2 is almost certainly right;
    only go to degree-3 if residuals show structure.
  \item \textbf{Saturation and dead-pixel masking.} Before any mean, mask
    pixels that are saturated ($\geq {\sim}1000$ DN on a 10-bit sensor) or
    identified as hot/dead from calibration. The IMX219 has ${\sim}0.1\%$
    bad pixels typically. Patch means should be trimmed means or medians
    over unmasked pixels.
  \item \textbf{Exposure bracketing.} I'd strongly consider capturing each
    LED condition at two or three exposures and selecting (or HDR-combining)
    to maximize SNR without saturating. The pale-yellow regime is where you
    need the dynamic range.
\end{itemize}

The output of this layer is: per-patch, per-LED, per-Bayer-channel, a
flat-fielded, dark-subtracted, outlier-masked pixel array. Not a mean yet
--- keep the pixels around, because the spatial distribution of values
inside the strip patch is itself information (see Layer~3).

\subsection*{Layer 2: Within-image normalization}

This is where your existing plan of ``divide by white patch'' lives, but I'd
do it differently than the original \texttt{.tex}.

\begin{itemize}
  \item After flat-fielding, the white reference patch should read
    approximately constant across the frame (that's what flat-fielding is
    for). So \emph{any} white patch on the card, or an average of several,
    gives you a single scalar per (LED, channel) to divide by. This is the
    illumination-intensity normalizer.
  \item The black patch, after dark subtraction, ideally reads zero. If it
    doesn't, you have stray light, and the black patch value is your
    stray-light estimate. Subtract it. (This is your stray-light test,
    built into every measurement.)
  \item The result is a \emph{reflectance-like} quantity per patch, per LED,
    per channel: dimensionless, between 0 and 1, interpretable as
    ``fraction of incident light reflected.''
\end{itemize}

This is a much cleaner foundation than the original paper's setup. It
cleanly separates intensity normalization (white) from stray-light
subtraction (black). Both are physical quantities, not statistical
nuisances.

\subsection*{Layer 3: Feature extraction from patch pixels}

Here I'd depart from the original paper significantly. The paper jumped
straight to ``patch mean, then Mahalanobis distance to references, with
pixel-noise covariance in the middle.'' I'd reconsider.

What you actually have from Layer~2 is, per strip, a multivariate
reflectance value per (LED, channel) position --- roughly $I\times 4$ values
if all Bayer channels are used for all LEDs, minus the ones knocked out by
low SNR. For typical cases that's maybe 8--12 numbers.

The useful features are:

\begin{itemize}
  \item \textbf{Mean reflectance vector.} The obvious one. Call it
    $\mathbf{r}\in\mathbb{R}^N$ where $N$ is the number of (LED, channel)
    combinations surviving SNR culling.
  \item \textbf{Within-patch variability.} The standard deviation or
    inter-quartile range of pixel reflectance inside the strip patch, per
    (LED, channel). This is a \emph{uniformity flag} --- if the strip
    reacted unevenly, these values will be elevated compared to
    calibration. Not a predictor of concentration, but a quality indicator.
  \item \textbf{Ratios / log-ratios between channels.} For a colorimetric
    measurement, ratios are often more informative than absolute values. A
    log-reflectance vector $-\log(\mathbf{r})$ is the right transform if
    you believe Beer-Lambert-like behavior (which is an approximation for
    strips, but a useful one).
\end{itemize}

I would \textbf{not} compute Mahalanobis distances to reference patches as
the primary feature. Here's why: reference patches are useful for
\emph{rebasing the measurement against known colors present in the same
image}, which corrects for illumination drift. But that's already handled
by white-patch normalization in Layer~2. The references give you a
\emph{check} --- does a ``yellow reference'' in this image read the same as
it did in calibration? --- but they aren't the input to the regression. The
regression input should be the strip's reflectance vector itself.

Keep the reference patches in the pipeline, but in a different role:
\textbf{as a per-image quality gate.} If any reference patch's reflectance
vector in the live image differs from its calibration mean by more than $k$
standard deviations (this \emph{is} where a Mahalanobis distance earns its
keep --- as an outlier detector on known-color patches), flag the image and
recalibrate or reject.

\subsection*{Layer 4: Dimensional reduction and covariate assembly}

You anticipated this: PCA on the strip reflectance vectors across a
calibration set will likely reveal that 1--2 components capture $>95\%$ of
the variance, because there's essentially one chemical axis
(white$\to$yellow$\to$brown). That's the ``principal color score'' $s_1$,
and possibly $s_2$.

The feature vector for regression is then small:
\begin{equation*}
  \mathbf{x} = \bigl(s_1,\; s_2,\; T_{\text{ambient}},\; T_{\text{water}},\;
    t_{\text{develop}},\; \text{strip-lot ID}\bigr).
\end{equation*}

Low-dimensional, interpretable, and each component has a physical meaning.
This is \emph{vastly} preferable to the original paper's $IJ+I$-dimensional
feature vector for a GP.

\subsection*{Layer 5: Regression}

Here the question is not ``GP vs.\ not GP'' but ``what's the simplest
model that works?'' In order of increasing complexity:

\begin{enumerate}
  \item \textbf{Univariate calibration curve.} $c = f(s_1)$. Fit a monotone
    spline, or a parametric curve motivated by the chemistry. This should
    be your baseline and it will probably cover 80\% of cases. It's
    intelligible to engineers. Every colorimetric analysis textbook uses
    this.
  \item \textbf{Bivariate.} $c = f(s_1, s_2)$. Include $s_2$ if residuals
    from (1) show structure along $s_2$.
  \item \textbf{With covariates.} Add $T_{\text{water}}$, $t_{\text{develop}}$,
    etc.\ if residuals correlate with them.
  \item \textbf{GP, only if needed.} Switch to a GP if a parametric model
    doesn't fit well and you have enough calibration data to justify it.
    GPs are powerful but opaque. Save them for when you've demonstrated
    that simpler models aren't sufficient.
\end{enumerate}

The original paper's ``GP with a 12+ dim feature vector'' is almost
certainly overfit for a calibration set of a few hundred strips.

\subsection*{Layer 6: Uncertainty calibration and decision threshold}

This is absent from the original and matters enormously for a field
instrument. Two things:

\begin{itemize}
  \item \textbf{Empirical prediction intervals via cross-validation or
    conformal prediction.} Rather than trust the model's internal
    uncertainty (GP posterior variance, or fitted spline confidence bands),
    compute residuals from leave-one-strip-out CV, and use their
    distribution to build honest 95\% prediction intervals. Report those.
  \item \textbf{Decision-relevant metrics.} The WHO limit is 10\,\textmu g/L.
    What matters is the \emph{false-negative rate at 10\,\textmu g/L}.
    Design the validation around this specifically. ROC-like analysis
    centered on the decision threshold is more important than overall RMSE.
\end{itemize}

% ============================================================
\section*{What I need to know to commit to specifics}
% ============================================================

Before I'd write down equations or code, I need the following, in roughly
this order of importance:

\begin{enumerate}
  \item \textbf{The actual strip chemistry.} Is this a Gutzeit
    (HgBr$_2$-on-paper) strip? Arsenator? Something else? The kinetics,
    temperature sensitivity, and spectral signature differ. This determines
    whether we expect Beer-Lambert-like log-reflectance behavior, what the
    dominant color axis looks like, and whether there's one reaction
    product or several.
  \item \textbf{A small real dataset.} Even 20--50 strips across the
    concentration range, imaged on the prototype, would let me see the
    actual shape of the color manifold, the noise level, and whether 1-D
    PCA really is enough. Without this we're just speculating.
  \item \textbf{The LED choices and their spectra.} You said ``blue, green,
    red, and optionally more.'' Which exact LEDs? Peak wavelengths, FWHMs?
    More LEDs aren't automatically better if their spectral bands don't
    align with features in the strip's reflectance spectrum.
  \item \textbf{The reference card design.} How many references, what
    colors? Are they printed (potentially fugitive), Munsell chips, or
    Pantone-painted? Long-term stability of the reference card is critical
    for an instrument that will live in the field.
  \item \textbf{Target performance.} LOD? Required precision at
    10\,\textmu g/L? Required precision at higher concentrations? Time
    budget per measurement? This drives how much calibration effort is
    justified.
  \item \textbf{Deployment constraints.} How much computation can a
    Raspberry Pi Zero do per measurement? A full GP inverse is fine for a
    few hundred points; a cross-validated conformal prediction interval
    means doing that several times. This bounds what's implementable.
\end{enumerate}

% ============================================================
\section*{What I recommend you learn or read}
% ============================================================

Not huge detours, but enough to let you judge proposals rather than accept
them:

\begin{itemize}
  \item \textbf{Beer-Lambert and diffuse reflectance basics.} Why
    log-reflectance is often the right transform. Kubelka-Munk is the
    standard model for reflectance from scattering media like paper
    strips; it's a one-page read and gives you a language for discussing
    what's going on optically. This is the single most useful thing to
    read.
  \item \textbf{PCA intuition, not mechanics.} You already have the
    mechanics from the \texttt{.tex}. What's worth building up is the
    intuition that PCA on a calibration set tells you the
    \emph{dimensionality of the chemistry}, and the explained-variance
    plot is a diagnostic tool you'll use to argue (to yourself and to
    reviewers) that a 1-D or 2-D model is justified.
  \item \textbf{Cross-validation and why ``training error'' lies.}
    Hastie, Tibshirani, Friedman (\emph{Elements of Statistical Learning})
    chapter~7 is the classic. The key idea: you will always be able to fit
    your calibration data well; the honest question is how well you predict
    \emph{unseen} strips, and CV is the answer.
  \item \textbf{Conformal prediction, at a high level.} This is a method
    for building honest prediction intervals around \emph{any} regression
    model without distributional assumptions. A short intro paper
    suffices. For a field instrument where the end user will make a
    health decision from the number, this matters.
  \item \textbf{What Mahalanobis distance actually is, and where it does
    and doesn't belong.} Not a deep dive --- the intuition ``distance
    measured in units of standard deviations, accounting for correlation
    between dimensions'' is enough, plus the fact that it requires you to
    specify \emph{which} covariance. In a classifier it's the within-class
    covariance. In outlier detection it's the expected-under-the-null
    covariance. In the original paper it was within-patch pixel noise,
    which was the wrong choice.
\end{itemize}

% ============================================================
\section*{What the rewritten document should look like}
% ============================================================

When we're ready to write, I'd structure it as:

\begin{enumerate}
  \item Instrument and measurement concept (hardware, chemistry, protocol,
    timing).
  \item Image processing pipeline (Layers 1--3 above) --- this is the
    engineering meat.
  \item Feature extraction and calibration model (Layers 4--5), presenting
    the univariate model as the baseline and more complex models as
    extensions justified by residual analysis.
  \item Uncertainty quantification and decision thresholds (Layer~6).
  \item Validation protocol and results (once data exists).
\end{enumerate}

Substantially shorter and substantially less mathematical than the
original. The original's weakness was that it paraded statistical
machinery; the new version's strength should be that each step has a
specific, defensible reason rooted in the physics and chemistry.

% ============================================================
\section*{Immediate next step I'd suggest}
% ============================================================

Get any real data --- even a quick prototype pass with ten strips at five
concentrations, imaged with whatever LEDs you have --- and look at it with
me. Everything above is informed speculation. A plot of $\log(1/r)$ vs.\
concentration for each LED, over a real run, answers more questions in ten
minutes than another round of theory.

\end{document}
